\cotefile{56}
\scanpage{15}
\diagram{\begin{tikzcd}[ampersand replacement=\&]
{} \& \deleted{\unsure{A^{+}|A}}\\
\deleted{\unsure{A|\mathcal A}}\arrow[ur] \& {}
\end{tikzcd}}

\[
\begin{aligned}
&\left\{\begin{aligned}
&\left.\begin{aligned}
H^i(E)&=0 \quad \deleted{\unsure{\text{soit}|i\ne0}}\quad
 \deleted{\text{Si }i\ne2(\nu-1)}\ 
 \inserted{i\ne0;\,2\unsure{(\nu-2)|\nu-2},\,4\nu-1}\\
H^{2(\nu-1)}(E)&\simeq\begin{cases}
0,\ l\ne 2\\
\mathbb Z/2\mathbb Z\quad \deleted{\text{si}}\ l=2
\end{cases}
\end{aligned}\right\}\\
&H^{4\nu-1}(E)\simeq\unsure{\mu|M}(-2\nu)
\end{aligned}\right.
&&\begin{aligned}\boxed{n+1=2\nu}\\ n=2\nu-1\end{aligned}
\end{aligned}
\]

\[
\begin{aligned}
&\left\{\begin{aligned}
H^{i}(E)&=0\qquad \text{si}\quad i\ne\inserted{0,}2\nu-1,2\nu+1,4\nu+1\\
H^0(E)&=\mu\unsure{(0)|w)|(w)}\\
H^{2\nu-1}(E)&\simeq\mu(-\nu+1)\\
H^{2\nu+1}(E)&\simeq\mu(-\nu-1)\\
H^{4\nu+1}(E)&\simeq\mu(-2\nu-1)
\end{aligned}\right.
&&\begin{aligned}\boxed{n+1=2\nu+1}\\n=2\nu\end{aligned}
\end{aligned}
\]
% archive sheet number: 14
\scanpage{16}
\mspage{\handwritten{16\unsure{a|.}}}
algébr\handwritten{iq}uement clos. Mais alors, $s$ désignant le point fermé de $S$, comme $G$ est étale et séparé sur $S$, il contient un sous-schéma à la fois ouvert et fermé $G'$, \inserted{fini sur $S$,} tel que $G'_s=G_s$ \inserted{(EGA II 6.2.6.)}, et on voit de suite que $G'$ est ici un sous-groupe de $G$. \deleted{\unsure{Par suite|Ensuite} la fibre géométrique $f(s)$} D'ailleurs $G'$ \deleted{\unsure{est|qui}} étant étale fini sur $S$ \deleted{\unsure{qui}} $=\operatorname{Spec}(V)$, $V$ complet à corps résiduel algébriquement clos, est un groupe constant, donc de la forme $A_S$, où $A\handwritten{=}G'(k(s))=G(k(s))$ a pour classe $f(s)$. \deleted{On voit donc} Si $B$ est la fibre géométrique de $G$ en le point générique $t$ de $S$, on a donc un \deleted{\unsure{isomor}}\inserted{monomor}phisme canonique $A\handwritten{\to}B$, ce qui prouve $f(s)\handwritten{\leq} f(t)$. (NB cette démonstration prouve en fait la semi-continuité pour une relation d'ordre sur $E$ plus fine que celle indiquée dans \handwritten{2}.1.). Le fait que $f$ soit continue en $s$ si et seulement si $s\handwritten{\rightsquigarrow}\operatorname{card} f(s)$ l'est résulte du fait que pour $w,w'\handwritten{\in}E$, les relations \deleted{\unsure{$w$|$w,$}} $w\handwritten{\leq}w'$ et $\operatorname{card} w=\operatorname{card} w'$ impliquent $w=w'$. Le fait que cette condition soit équivalente à\deleted{\unsure{s|,}} la finitude de $W$ sur un voisinage de $s$ est alors indépendant de\deleted{\unsure{s}} \inserted{la} structure\deleted{\unsure{s}} de groupe\deleted{\unsure{s}} \inserted{\handwritten{sur $W$}}, et a été signalé après \inserted{\handwritten{Exp $\overline{\underline{\mathrm{XI}}}$}} 5.10. ; d'ailleurs sa démonstration \deleted{\unsure{par}} se fait aisément par les arguments précédents, \inserted{en} utilisant le c\handwritten{ri}tère valuatif de propreté EGA II 7.3.8.

\uline{Lemme} \handwritten{2}.3. Soient $G$ un groupe algébrique affine lisse sur un corps algébriquement clos $k$, $R\handwritten{\subset}T$ deux sous-tores, $W(R)$ et $W(T)$ les deux groupes finis associés comme dans \inserted{\handwritten{Exp $\overline{\underline{\mathrm{XI}}}$,}} 5.9. , quotient\handwritten{s} du normalisateu\unsure{r} par le centralisateur. Alors $W(R)$ est isomorphe à un quotient d'un sous-groupe de $W(T)$.

\deleted{Dans la démonstration qui suit, il sera commode d'interpréter $G,R,T,C(R),N(R),C(T),N(T)$ de façon purement ensembliste, comme les groupes de points des préschémas $G,R,T$ dans $k$, et les centralisateurs et nor}
% archive sheet number: 15
\scanpage{17}
\uline{Cohomologie à supports compacts} \deleted{\unsure{de}} de $C,E,\widehat E$ ....

La cohomologie à \unsure{support compact|supports compacts} des $E,\widehat E$ se déduit simplement par dualité des \deleted{celles} \inserted{\unsure{la Ch. usuelle|les Ch. usuelles}} des $E,\widehat E$ lorsque $X$ est lisse ... Celle de $\widehat C$ est identique à sa cohomologie usuelle (car $\widehat C$ est propre...).

Reste à déterminer la ch. à supports compacts de $C$.

\unsure{Car.|Or.}
\[
\deleted{H_c^*(C')\longleftarrow H_c^{\unsure{p|*}}(\unsure{\cdot|\ldots})}
\]

\diagram{\begin{tikzcd}[ampersand replacement=\&,column sep=tiny,row sep=small]
\cdots\arrow[r,dash] \& H^i(\widehat C)\arrow[r,dash]\arrow[d,phantom,"\rotatebox{90}{$\simeq$}" description] \& H^i(X)\arrow[r,dash] \& H_c^{i+1}(C)\arrow[r,dash] \& H^{i+1}(\widehat C)\arrow[r,dash]\arrow[d,phantom,"\rotatebox{90}{$\simeq$}" description] \& H^{i+1}(X)\\
{} \& \substack{H^{i-2}(X)(-1)\\ \text{si }i\ne0} \& {} \& {} \& H^{i-1}(X)(-1) \& {}
\end{tikzcd}}

Comparant avec la suite exacte analogue pour déterminer \inserted{la ch. de} $E$ \inserted{\unsure{ou celle donnant} $H_x^*(\widehat C)$ \unsure{pr|par} $\widehat E\subset\widehat C$,} on est amené à prévoir $H^i(E)\overset{\text{si }i\ne0}{\simeq}H_c^{i+1}(C)\mathrel{\underset{\unsure{a|\sim}}{\simeq}}H_{\deleted{\unsure{c|x}}x}^{i+1}(C)$, \deleted{\unsure{et}} \deleted{\unsure{on a en effet}, \unsure{grâce à l'inclusion} $E\subset C$,}
\[
\boxed{H_x^*(C)\simeq H_c^*(C)\quad ??}
\]
\[
\deleted{H_x^i(C)=H_c^i(C)=}
\]
% archive sheet number: 16
\scanpage{18}
\mspage{\handwritten{\deleted{\unsure{48|18}}4}}
\handwritten{l'}exposé \deleted{ultérieur} \inserted{\handwritten{$\overline{\underline{\mathrm{XIV}}}$}} \deleted{que la réciproque est vraie} \inserted{\handwritten{que la réciproque est vraie}} lorsque $S$ est \deleted{\unsure{le spectre}} \inserted{artinien}\handwritten{,} \deleted{\unsure{d'un corps},} \deleted{\unsure{d'un ann. artin|anneau artin}} \handwritten{ou} \deleted{\unsure{ou}} lorsque $S$ est un schéma local et $G$ est \handwritten{\textquotedblleft}réductif'' : dans ce cas, tout tore de $G$ est contenu dans un tore m\handwritten{a}ximal.

\uline{Définition} \handwritten{1}.5. Soit $G$ un groupe algébrique sur un corps $k$\handwritten{\unsure{.|,}} \deleted{\unsure{et soit}} On appelle rang réductif (resp. rang nilpotent, resp. rang unipotent ..) de $G$ le rang réductif (resp. ...) de $G_{\bar k}$, où $\bar k$ est une clôture algébrique de $k$.

On notera qu'en vertu de \handwritten{1}.2. et de la commutation de la formation de $\underline{\unsure{\mathbb C|C}\mathrm{entr}}_G(T)$ avec l'extension de la base, \deleted{\unsure{que}} les \deleted{rangs} notions de rang introduites dans \handwritten{1}.5. sont \uline{invariantes par extension du corps de base}; d'autre part, pour $k$ algébriquement clos, elles coïncident avec celles introduites au début du présent numéro.

\uline{Remarque} \handwritten{1}.6. Il n'est pas difficile de construire un schéma en groupes affine et lisse $G$ sur \deleted{\unsure{n}}le spectre $S$ d'un anneau de valuation discrète, dont la fibre générale soit isomorphe à $\underline G_m$, et la \deleted{\unsure{s}}fibre spéciale isomorphe à $\underline G_a$. Un tel $G$ ne contient aucun tore sauf le tore trivial \inserted{$T$} (réduit au sous-groupe unité), qui n'est évidemment pas tore maximal. De façon précise, dans la fibre spéciale $\unsure{\mathbb G|G}_{\unsure{o|0}}=\underline G_{a\ k}$, $T_o$ est bien tore maximal, mais dans la fibre générale $G_1=\underline G_{m\ K}$, $T_1$ n'est plus tore maximal (NB $k$ est le corps résiduel, $K$ le corps des fractions de l'anneau de valuation envisagé). On voit aussi sur cet exemple que le rang réductif de $G_s$ ($s\handwritten{\in}S$) n'est pas fonction continue de $s$. On a cependant les résultats suivants:

\uline{Théorème} \handwritten{1}.7. Soi\handwritten{e}nt $S$ un préschéma, $G$ un $S$-préschéma en groupes affine et lisse sur $S$. Pour tout $s\handwritten{\in}S$, considérons $\handwritten{\rho}_r(s)=\handwritten{\rho}_r(G_s)$ et $\handwritten{\rho}_n(s)=\handwritten{\rho}_n(G_s)$, les rangs réductif\deleted{\unsure{s}} et nilpotent de $G_s$ (\handwritten{1}.5.). Avec ces notations, on a ce qui suit:
% archive sheet number: 17
\scanpage{20}
\uline{Cohomologie des quadriques}

$Q^n$ quadrique non singulière de dim $n$ sur $k$ sép. clos.

$Q^{n-1}$ section à l'infini.

\[
U^n=AQ^n=Q^n-Q^{n-1}
\qquad\text{\unsure{supposons|pposons|pour} $A$ coefficient annulé par $m$}
\]
\[
\begin{aligned}
H_c^i(U^n,A)&\simeq
\begin{cases}
0&\text{si }i\ne n,2n\\
A\otimes V'_n\deleted{\unsure{(i)|0}}&\text{si }i=n\\
A\otimes\mu_m(-n)&\text{si }i=2n
\end{cases}\quad\inserted{\text{si }n\ne0}\\
V'_n&=\check V_n\otimes\mu_m(-n)=\begin{cases}
\mu_m(-\frac n2)&n\text{ pair}\\
\mu_m(\frac{\unsure{-(n+1)|-n+1}}2)&n\text{ imp.}
\end{cases}\\
V'_n&=V_n^{\unsure{*|+|\vee}}=\mu_m(-\tfrac n2)\qquad n\text{ pair}\\
V'_n&\unsure{\simeq|=} V_n\otimes\mu_m(1)\qquad n\text{ impair}
\end{aligned}
\]
% slanted
\margin{Si $n$ impair, ces isom. ne sont déterminés qu'au signe près ....}
\[
\Updownarrow
\]
dualement
\[
\begin{aligned}
H^i(U^n,A)&\simeq\begin{cases}
0 &\text{si }i\ne0,n\\
A\otimes V_n &\text{si }i=0\\
A &\text{si }i=n.
\end{cases}\quad\inserted{\text{si }n\ne0\ \substack{\nearrow\\\nearrow}}\\
V_n&\simeq\mu_m(-\tfrac n2)\quad\text{si $n$ pair,}\\
V_n&\simeq\mu_m(\tfrac{\unsure{-n+1|-(n+1)|n\pm1}}2)\quad\text{si $n$ impair}
\end{aligned}
\]
$H_c^i(U^n,A)\xrightarrow{\varphi}H_c^i(U^n,A)$ est \uline{nul} \uline{sauf} si $n$ pair, $i=n$, dans ce cas \unsure{devient} un \unsure{isomorphisme} \inserted{exc. 2}.
\[
\deleted{\unsure{H^i(Q^n)\simeq H^i(Q^n)|H^i(Q)\simeq H^i(Q)}\longrightarrow}
\]
Il y a \inserted{deux} \unsure{un gén.|un génér.} $l$ de $H_c^i(U^n,A)$ tels que $\langle l,\varphi(l)\rangle=2$.

\diagram{\begin{tikzcd}[ampersand replacement=\&,column sep=tiny]
H^{i-1}(U^n)\arrow[r,dash] \& H^{i-2}(Q^{n-1})(-1)\arrow[r] \& H^i(Q^n)\arrow[r] \& H^i(\deleted{\unsure{Q}}U^n)\arrow[r,dash] \& H^{i-1}(Q^{n-1})(-1)\arrow[r,dash] \& \cdots
\end{tikzcd}}
\[
H^{i-2}(Q^{n-1})(-1)\xrightarrow{\alpha_*}H^i(Q^n)\quad\text{est}\quad
\left|\begin{aligned}
&\text{surjectif si }i\ne0,n\\
&\text{injectif si }i\ne1,n+1
\end{aligned}\right.
\quad\left|\begin{aligned}&\text{bijectif si}\\&i\ne0,\deleted{\unsure{1}}n,n+1\end{aligned}\right.
\]

Que se passe-t-il \unsure{pr|pour} $i=0,1,n,n+1$ ??

Pour $i=0,1$ $\alpha_*$ est \uline{nul} [$i=1$, \deleted{bij} $\alpha$ bijectif, car les deux mbs nuls)

Pour $i=n,n+1$ on a la \deleted{\unsure{suite}} \unsure{exacte} (si $\boxed{n\ne1}$ \unsure{donc} $n-1\ne0$]
\diagram{\begin{tikzcd}[ampersand replacement=\&,column sep=tiny]
0\arrow[r] \& H^{n-2}(Q^{n-1})\deleted{(-1)}\arrow[r] \& H^n(\deleted{\unsure{U}}Q^n)\arrow[r] \& H^n(U^n)\arrow[r,dash] \& H^{n-1}(Q^{n-1})(-1)\arrow[r,dash] \& H^{n+1}(Q^n)\arrow[r,dash] \& 0
\end{tikzcd}}

1) $n$ \uline{pair} $>0$, la suite exacte se réduit à
\diagram{\begin{tikzcd}[ampersand replacement=\&,column sep=small]
0\arrow[r] \& H^{n-2}(Q^{n-1})\deleted{(-1)}\arrow[r] \& H^n(Q^n)\arrow[r] \& H^n(U^n)\arrow[r] \& 0
\end{tikzcd}}
donc $H^n(Q^n)$ est extension de $H^n(U^n)$ par $H^{n-2}(Q^{n-1})(-1)$ ...
% archive sheet number: 19
\scanpage{21}
\mspage{\handwritten{\unsure{15|35}}}
me $R_s$ est un tore maximal dans $G_s$, et la formation de \uline{Norm}, \uline{Centr}, \uline{Norm}/\uline{Centr} est compatible avec toute extension de la base, en particulier le passage aux fibr\unsure{e|è|é}s, on voit que l'on a
\[
w(s)=w'(s)\qquad ;
\]
je dis que de plus, pour $t$ voisin de $s$, on aura les inégalités
\[
w(t)\handwritten{\geq}w'(t)\qquad\text{et}\qquad w'(t)\handwritten{\geq}w'(s)\quad ,
\]
ce qui suffira pour établir a). \deleted{La deuxième} \handwritten{\unsure{C|L}}es deux inégalités \deleted{\unsure{résulte|résultera}} sont contenu\handwritten{es} dans les deux lemmes suivants:

\uline{Lemme} \handwritten{2}.2. Soi\deleted{\unsure{t}}ent $S$ un préschéma, $W$ un $S$-préschéma en groupes qui est étale, \deleted{quasi-fini et} séparé et quasi-fini sur $S$. Pour tout $s\handwritten{\in}S$, soit $f(s)$ la classe de la fibre géométrique de $W$ en $s$, qui est un \deleted{ensemble}\inserted{élément} de l'ensemble ordonné $E$ des classes de groupes finis à isomorphisme près, introduit dans \handwritten{2}.1. Alors la fonction $f:S\handwritten{\to}E$ est semi-continue \deleted{\unsure{sup}}inférieurement. Pour qu'elle soit constante au voisinage de $s\handwritten{\in}S$\handwritten{,} \deleted{(ce qui revient au même, pour} i\deleted{\unsure{o}}l faut et suffit qu'il en soit de même de la fonction $s\handwritten{\rightsquigarrow}\operatorname{card}f(s)$, et pour ceci il faut et suffit que \handwritten{$W$} soit fini sur $S$ au dessus d'un voisinage ouvert $U$ de $s$.

Ce résultat est un raffinement, en termes de groupes, de celui invoqué dans la démonstration de \inserted{\handwritten{Exp $\overline{\underline{\mathrm{XI}}}$}} \handwritten{5}.10. , nous nous bornons à une esquisse de \inserted{\handwritten{la}} démonstration (\inserted{\handwritten{qui est}} d\handwritten{u} type \handwritten{le} plus standart). On est ramené comme d'habitude au cas $S$ affine noethérien. On voit de suite que la fo\deleted{\unsure{u}}nction $f$ est constructible (EGA \deleted{\unsure{IV|II}} $\unsure{0|O}_{\mathrm{III}}$ 9.3.1. et 9.3.2. , et sor\handwritten{it}es de EGA IV 9), et en vertu de EGA $0_{\mathrm{III}}$ 9.3.4. on est ramené \inserted{pour la semi-continuité} à prouver que si $t$ est une générisation de $s$, on a $f(t)\handwritten{\geq}f(\deleted{\unsure{x}}s)$. Cela nous ramène \deleted{au cas} grâce à EGA II 7.1.7. au cas où $S$ est le spectre d'un anneau de valuation discrète, qu'on peut supposer complet à corps résiduel
% archive sheet number: 20
\scanpage{22}
\unsure{b|2}) $n$ impair $\geq1$, la suite exacte se \unsure{lit alors|réduit} :
\diagram{\begin{tikzcd}[ampersand replacement=\&,column sep=small]
0\arrow[r] \& H^n(U^n)\arrow[r] \& H^{n-1}(Q^{\unsure{n+1|n-1}})(-1)\arrow[r,dash] \& H^{n+1}(Q^n)\arrow[r,dash] \& 0
\end{tikzcd}}
i.e. $H^{n\unsure{+|-}1}(Q^{n-1})$ est une extension de $H^{n+1}(Q^n)$ par

a) Cohomologie de $Q^n$ nulle en dim impaire.

b) $H^i(U^n)$ nul si $i\ne0,n$, isomorphe à $A$ (resp. $A\otimes \unsure{V'|V|\nu}_n$) si \inserted{$i=0$ (resp. $i=n$)}

b') Énoncé dual

c) $H^{i-2}(Q^{n-1})\unsure{\deleted{(-1)}|(-1)|\deleted{(A)}}\longrightarrow H^{\deleted{\unsure{2}}i}(Q^n)$ est
\[
\left|\begin{aligned}
&\left.\begin{aligned}
&\text{bijectif si }\deleted{\unsure{i\ne0,\ 2\leq i}}\ \inserted{i\geq2},\ i\ne n\\
&\text{injectif pour tt }i
\end{aligned}\right\}\ (n\ \text{\uline{pair}})\\
&\left.\begin{aligned}
&\text{bijectif si }i\geq2,\ i\ne n+1\\
&\text{surjectif pour tt }i\geq1
\end{aligned}\right\}\ (n\text{ impair})
\end{aligned}\right.
\]

c') Énoncé dual : $H^i(Q^n)\longrightarrow H^i(Q^{n-1})$ est
\[
\left|\begin{aligned}
&\left.\begin{aligned}
&\text{bijectif si }i\leq 2n-2,\ i\ne n\\
&\text{\unsure{surjectif|injectif} pour tt } \unsure{i|\text{''}}
\end{aligned}\right\}\ \text{\uline{$n$ pair}}\\
&\left.\begin{aligned}
&\text{bijectif si }i\leq2n-1,\ i\ne n-1\\
&\text{injectif pour tt }i\leq2n-1
\end{aligned}\right\}\ n\text{ impair}
\end{aligned}\right.
\]

d) $H^{\unsure{i|2i}}(Q^n)$ \deleted{de rang} \deleted{\unsure{en|est}} de rang 1 en \deleted{\unsure{toute}} dimension \inserted{paire} \uline{sauf} pour $n$ pair et \deleted{\unsure{si}} $i=n$, où il est de rang 2.

e) Si $n$ pair,

% margin note written sideways
% Inverted pencil group, including its marginal F, read after a half turn.
\diagram{\begin{tikzcd}[ampersand replacement=\&,column sep=small]
F \& {}\\
X\arrow[r,phantom,"\subset" description]\arrow[dr,dash] \& X'=Y[\unsure{h|h_1|t_1},-,\unsure{h_n|h|t_n}]\arrow[d,dash]\\
{} \& Y
\end{tikzcd}}
% archive sheet number: 21
\scanpage{23}
\mspage{\handwritten{23}}
multiplicatif \deleted{\unsure{autre}} sauf le groupe unité. Comme un sous-groupe d'un \margin{$\bullet$} groupe de type multiplicatif est de type multiplicatif (Exp VI\quad ), il s'ensuit que \deleted{tout} pour tout sous-groupe de type multiplicatif $H$ de $G/Z_{\unsure{o|0}}$, on a $H\handwritten{\cap}G/Z_o=0$, donc $H\handwritten{\to}G/\deleted{\unsure{Z}}G_o$ est injectif, Comme $G/G_o$ est un groupe algébrique radiciel, donc fini sur $k$, cela implique que $H$ est lui-même radiciel et de \uline{rang ma\deleted{\unsure{h}}joré} par celui de $G/G_o$. Cela \handwritten{impl\unsure{\inserted{i}q|iq}ue} que la famille des sous-groupes de type multiplicatifs de $G$ admet un plus grand élément, soit $Z$.

Je dis que $G/Z$ n'a pas de sous-groupe de type multiplicatif autre que le sous-groupe unité. Cela résulte du fait (démontré dans \margin{$\bullet$} Exp VII\quad ) qu\deleted{e}'une extension commutative de deux groupes algébriques de type multiplicatif est de type multiplicatif\handwritten{;} \deleted{compte tenu du caractè-} \handwritten{si} donc $Z'$ est l'image inverse dans $G$ d'un sous-groupe de type multiplicatif de $G/Z$, alors $Z'$ est de type multiplicatif, donc $Z'=Z$\unsure{c} par le caractère maximal de $Z$, donc $Z'/Z=0$.

Il résulte maintenant aisément du \textquotedblleft principe de l'extension finie'' que pour toute extension $K$ de $k$, $(G/Z)_K=G_K/Z_K$ n'a pas non plus de sous-groupe de type multiplicatif sauf le groupe unité.

Nous pouvons maintenant prouver que $Z$ est un centre réductif de $G$. En effet, soit $u:H\handwritten{\to}G_S$ un homomorphisme de $S$-groupes, où $S$ est un préschéma sur $k$ et $\unsure{\mathbb H|H}$ un groupe de type multiplicatif et de type fini sur $S$, prouvons que $u$ se factorise à travers $\unsure{\mathbb Z|Z}_S$. Il revient au même de dire que l'homomorphisme composé $H\handwritten{\to}G_S\handwritten{\longrightarrow}(G/Z)_S=G_S/Z_S$ est nul. Or je dis que, posant $U=G/Z$, tout homomorphisme $u:H\handwritten{\to}U_S$ est nul. En effet, en vertu de Exp IX 5.2. on est ramené au cas où $S$ est le spectre d'un corps, et par Exp IX 6.8. cela résulte alors du
% archive sheet number: 22
\scanpage{24}
\[
\unsure{\mathbb P|P}^1\times\mathbb P^1\simeq Q^2\subset Q^1\simeq P^1,\qquad
\frac{n-1}{2}+1=\frac{n+1}{2}
\]
\diagram{\begin{tikzcd}[ampersand replacement=\&]
\cdot\arrow[r,dash] \& H^0(Q^1)\arrow[r] \& H^2(\\
{} \& \unsure{\checkmark} \& {}
\end{tikzcd}}
\[
\begin{aligned}
&& n=2\nu\\
H^2(Q^n)&\mathrel{\rule{1.6em}{.35pt}}& H^{2\nu-2}(Q^n)
\end{aligned}
\]
\[
\begin{aligned}
\xi^{0}&\ \mathord{-}\ \xi^{\nu-1}&\xi^\nu&\quad \deleted{\unsure{\xi}}\xi^{\nu+1\,\unsure{[\tfrac12|[2|/2|[\xi}}\ \cdots\ \xi^{2\nu}\\
&&\unsure{l|h|\ell}&\\
\xi^{\unsure{1|4}}_{2}&\quad \deleted{\cdots\ \xi_{\nu}\qquad\xi}&&
\end{aligned}
\]
\[
Q^n\qquad\left\{\begin{aligned}
&\xi_i\in H^{2i}(Q^n,\unsure{\mu|M}(i))\qquad\deleted{0\leq i\ \unsure{\text{pair}}\ \leq2n}\\
&\deleted{\lambda_{\unsure{2\nu|\nu}}\in H^{2\nu}(Q^{2\nu})}\quad\lambda_{2\nu}\in H^{2\nu}(Q^{2\nu})
\end{aligned}\right.
\]
\[
\deleted{Q}\qquad \xi_{\unsure{i|0}}
\]
\diagram{\begin{tikzcd}[ampersand replacement=\&,column sep=small]
Q^{n-1}\arrow[r,phantom,"\subset" description] \& Q^n\\
\xi_{\unsure{i|1}} \& \xi_i\arrow[l]\\
\xi_i\arrow[r,decorate,decoration={snake,amplitude=.26mm,segment length=1.7mm,post length=1.4mm}] \& \xi_{i+2}
\end{tikzcd}}
\[
\begin{pmatrix}1,1\\1,-1\end{pmatrix}
\]
\diagram{$\begin{array}{c@{\;}c@{\;}c@{\;}c@{\;}c@{\;}c@{\;}c@{\;}c@{\;}c}
Q^0&\subset&Q^1&\subset&Q^2&\subset&Q^3&\subset&Q^4\\
\lambda_0^{\unsure{*|0|+}},\xi^0_{\unsure{i|0|+}}&&\xi^0&&\xi^0&&\xi^0&&\xi^0\\
 &&\frac12\xi&&(\lambda_2,\xi)&&\xi&&\xi\\
 &&&&\frac12\xi^2&&\frac12\xi^2&&(\lambda_{\deleted{\unsure{2}}4},\xi^2)\\
 &&&&&&\frac12\xi^3&&\frac12\xi^3\\
 &&&&&&&&\frac12\xi^4
\end{array}$}
\[
\deleted{Q^{\unsure{n=2\nu}}}:\quad \deleted{\xi}\ \cdots
\]
% slanted
\margin{d'indice 2 dans $H^{2\nu}(Q^{2\nu})$, il faut \unsure{rajouter|ajouter} $\frac12(\xi^\nu+\lambda_\nu)=\varphi_\nu$ pour avoir une base $(\xi^\nu,\varphi_\nu)$}
\[
\boxed{\begin{aligned}
Q^{n=2\nu}:\quad &\xi^0\ \cdots,\xi^{\nu-1},[\xi^\nu,\lambda_\nu],\tfrac12\xi^{\nu+1},\cdots,\tfrac12\xi^{2\nu}\\
Q^{n=2\nu+1}:\quad &\xi^0,\cdots,\xi^\nu,\tfrac12\xi^{\nu+1},\cdots,\tfrac12\xi^{2\nu+1}
\end{aligned}}\qquad \lambda_\nu
\]
% archive sheet number: 23
\scanpage{25}
\mspage{\handwritten{\unsure{11|15}\deleted{\unsure{a}}}}
\deleted{\unsure{G de R.|de R.} En vertu de 10.5.3. c'est un sous-groupe fermé de $G$ lisse sur $S$, donc lisse et affine sur $S$, et comme $T$ est central, on a $T\unsure{\subset|=}H$.}

Ceci résulte du \deleted{résultat suivant}

\uline{Lemme} \handwritten{1}.14. Soient $G$ \deleted{affine}\inserted{\handwritten{un $S$-préschéma en groupes,}} de présentation finie sur $S$, $T$ un tore maximal de $G$, \deleted{alors $T$ contient tout sous-tore \uline{central}} $R$ un sous-tore de $G$, et supposons que $T$ et $R$ commutent. Alors $R\handwritten{\subset} T$.

En effet, comme $R$ et $T$ commutent, le morphisme $R\times_S T\handwritten{\to}G$ défini \deleted{\unsure{pa|par}} par $(r,t)\handwritten{\rightsquigarrow}r.t$ est un homomorphisme de groupes, donc en vertu de \handwritten{1}.8. il admet un sous-groupe image $T'$ dans $G$, qui est un \deleted{sous-}groupe de type multiplica\unsure{ti|it}f quotient de $R\times_S T$, donc un tore, contenant évidemment $T$. Comme $T$ est un tore maximal, on a $T'=T$ (\handwritten{1}.4.), donc $R\handwritten{\subset}T$, \deleted{cqfd.} ce qui démontre \handwritten{1}.14. donc \handwritten{1}.13. En particulier\handwritten{, utilisant 6.7. b):}

\uline{Corollaire} \handwritten{1}.15. Soit $G$ un $S$-préschéma en groupes \deleted{\unsure{affi|affine}} \uline{commutatif}, lisse et affine sur $S$\handwritten{,} \inserted{\handwritten{de rang réductif localement constant.}} Alors $G$ admet u\deleted{\unsure{j}}n unique tore maximal, et c\deleted{\unsure{l}}e\deleted{\unsure{lui|t}} \deleted{qui est le plus} dernier contient tout sous-tore de $G$.

\uline{Corollaire} \handwritten{1}.16. Soit $G$ un $S$-préschéma en groupes, lisse et affine sur $S$. Supposons que pour tout $s\handwritten{\in}S$, désignant par $\bar s$ le spectre d'une clôture algébrique \inserted{$k$} de $k(s)$, la fibre géométrique $G_{\bar s}$ soit un groupe algébrique connexe, et nilpotent (au sens de BIBLE, i.e. le groupe $G(k)$ de ses points à va\deleted{\unsure{b}}leurs dans $k$ est nilpotent). \inserted{Supposons de plus le rang réductif de $G$ loc.constant.} Alors $G$ admet u\deleted{\unsure{j}}n unique tore maximal $T$, \inserted{\handwritten{de plus}} $T$ est central et c'est le plus grand sous-tore de $G$.

En vertu de \handwritten{1}.7. b), $G$ admet localement pour fpqc un tore maximal, et en vertu de \handwritten{1}.13. on est ramené à prouver qu'un tore maximal de $G$ est central. En vertu de Exp \handwritten{$\overline{\underline{\mathrm{IX}}}$}\deleted{\illegible{1}}, 5.6. b), on est ramené au cas où $S$ est le spectre d'un corps, qu'on peut supposer algébriquement clos.
% archive sheet number: 24
\scanpage{26}
\[
Q^i\subset Q^j \qquad\qquad \xi
\]
\[
\deleted{\mathbb Z[T][T']/T^{2\nu+2},\ 2T'-T^{\nu+1}}
\]
\[
\left\{\begin{aligned}
&\xi^0,\cdots\,\xi^{\nu-1},[\xi^\nu,\lambda_\nu],\ \xi_{\nu+1},\cdots,\xi_{2\nu}\\
&\xi^0,\cdots,\xi^\nu,\ \xi_{\nu+1},\cdots,\xi_{2\nu+1}
\end{aligned}\right.
\]
\[
\left[\begin{aligned}
2\xi_i&=\xi^i\\
\xi^{i}_{\deleted{\unsure{i|2}}}\xi_j&=\xi_{i+j}\\
\xi_i\xi_j&=0\quad\text{pour raison de degré}\\
\lambda_\nu^2&=2\xi_{2\nu}=\xi^{2\nu}\\
\lambda_\nu\xi&=0\\
\lambda_\nu\xi_\alpha&=0\quad\alpha\geq\nu+1\\[-2pt]
&\text{raisons de degré}
\end{aligned}\right.
\]
\[
Q^i\subset Q^j\qquad (i\leq j)
\]
\diagram{\begin{tikzcd}[ampersand replacement=\&,column sep=small,row sep=small]
\xi^\alpha \& \xi^\alpha\arrow[l,decorate,decoration={snake,amplitude=.26mm,segment length=1.7mm,post length=1.4mm},"\deleted{\unsure{\xi|\xi_\alpha|\nearrow}}" description]\\
\xi_\alpha \& \xi_\alpha\arrow[l,decorate,decoration={snake,amplitude=.26mm,segment length=1.7mm,post length=1.4mm}]\\
0 \& \lambda_\nu\arrow[l,decorate,decoration={snake,amplitude=.26mm,segment length=1.7mm,post length=1.4mm}] \& {}\arrow[d,phantom,"{\Big\}\ \text{si }i<j=2\nu}" description]\\
\xi_\nu \& \varphi_\nu\arrow[l,decorate,decoration={snake,amplitude=.26mm,segment length=1.7mm,post length=1.4mm}] \& {}\\
\xi^\alpha\arrow[r,decorate,decoration={snake,amplitude=.26mm,segment length=1.7mm,post length=1.4mm}] \& \xi^{\alpha+(j-i)} \& {}\\
\xi_\alpha\arrow[r,decorate,decoration={snake,amplitude=.26mm,segment length=1.7mm,post length=1.4mm}] \& \xi_{\alpha+(j-i)} \& {}\\
\lambda_\nu\arrow[r,decorate,decoration={snake,amplitude=.26mm,segment length=1.7mm,post length=1.4mm}] \& 0 \& {}\arrow[d,phantom,"{\Big\}\ \text{si }i=2\nu<j}" description]\\
\varphi_\nu\arrow[r,decorate,decoration={snake,amplitude=.26mm,segment length=1.7mm,post length=1.4mm}] \& \xi_{\nu+(j-i)} \& {}
\end{tikzcd}}
\margin{: $\lambda_\nu$ est une \uline{classe \unsure{effective|ineffective}}}
\margin{: $\lambda_\nu$ est une classe \textquotedblleft évanescente''.}
% archive sheet number: 25
\scanpage{27}
\mspage{\handwritten{21}}
finie, \inserted{\handwritten{et à fibres affines,}} $Z$ un sous-préschéma en groupes. Pour que $Z$ soit un centre réductif de $G$, il faut et il suffit qu'il soit de type multiplicatif, et que pour tout $s\handwritten{\in}S$, \deleted{désignant par} $Z_s$ soit \deleted{\unsure{le}}un centre réductif de $G_s$.

En effet, il résulte d'abord de Exp IX 5.6. b\deleted{\unsure{1}}) qu'alors $Z$ est central. Comme les conditions envisagées sont stables par changement de base, il reste à prouver que tout homomorphisme central $u:H\handwritten{\to}G$, avec $H$ de type multiplicatif et de type fini, se factorise par $Z$. Or , $Z$ étant central, $u$ et l'immersion canonique $v:Z\handwritten{\to}G$ définissent un homomorphisme de groupes
\[
w:H\times_S Z\handwritten{\longrightarrow}G\handwritten{.}
\]
\handwritten{E}n vertu de Exp IX 6.8. ce dernier admet un groupe image $K$, qui est un sous-groupe de type multiplicatif de $G$, et tout revient à prouver que l'on a $K=Z$. Or cela est vrai fibre par fibre d'après l'hypothèse sur $Z$, et il suffit maintenant d'appliquer Exp IX 5.1. bi\handwritten{s,} \inserted{ce qui achève de prouver \handwritten{4}.3.}

On notera que dans le critère \handwritten{4}.3., l'hypothèse que pour tout $s\handwritten{\in}S$, $Z_s$ soit \deleted{\unsure{le}}un centre réductif de $G_s$ est en fait purement géométrique, i.e. il suffit de la vérifier sur la clôture algébrique de $k(s)$, comme il résulte de la deuxième assertion \handwritten{4}.2.

\margin{\handwritten{\unsure{$\nearrow$|1}}}
\uline{Théorème} \handwritten{4}.4. Soit $G$ un groupe algébrique \uline{affine} sur un corps $k$. Alors $G$ admet un centre réductif.

Comme le centre de $G$ est représentable par un sous-groupe fermé de $G$ (Exp VIII 6.7.) , on est ramené aussitôt au cas où $G$ est commuta\deleted{\unsure{ti|tif}}tif. De plus, en vertu de \handwritten{4}.2. on peut supposer $k$ algébriquement clos. Nous verrons \deleted{\unsure{plus loin|plus bas}}\inserted{\handwritten{dans Exp $\overline{\underline{\mathrm{XIII}}}$}} que dans ce cas $G$ s'écrit comme un produit \deleted{\unsure{et|de}} $Z\times U$, où $Z$ est de type multiplicatif, et $U$ \textquotedblleft unipotent'', d'où il résultera aussitôt que $Z$ est bien un centre réductif de $G$. Nous allons
% archive sheet number: 26
